% Tapis de Quillen, with Quillen's offprint and preprint (cote 162-5, typed part): archive cote 162-5, site page 9 (PDF page 10).
% typescript. Draft compared: none.
% AI-transcribed by Claude Opus 5.5 (claude-opus-5-5, adaptive thinking, effort high), Message Batches API, batch t162_full (msgbatch_01UAh3zyWPd454AeDkhnmQ33), 2026-11-13; not yet checked by eye.
\scanpage{162-5/9}
\mspage{\hand{8}}

% NOTE: the right edge of the typescript is cut off on the scan; line ends are transcribed as far as visible
On doit trouver une approximation meilleure d'un cran à la classifi
tion topologique en prenant des  (2-Cat)-groupes resp. des (2-Cat)-group
abéliens,  ou plus généralement des " Gr-2-catégories "  resp.
"\unsure{Grab}-2-catégories" . (NB  \unsure{Une} Grab-catégorie défi\xout{\unsure{k}}nit une Grab-2catégor
comme une catégorie définit une 2-catégorie; du point de vue semi-simpli
al, cette opération serait analogue à celle de "cosquélette$^{\text{x}}$ tronqué";
Mais également une Grab-2-catégorie défini\over{e}{t} \over{en}{une} G\unsure{r}ab-catégorie$^{\text{x}}$; savoir la
catégorie des automorphismes de l'objet unité; cette \xout{\illegible{1}} opération serait
analogue à celle de "tuage des groupes d'homotopie en dim \hand{$>$}\unsure{\over{2}{1}}"). Pour l
notion gro\xout{i}ssière ( Gra\xout{b} ou Grab) l'idée de ma lettre à Verdier était qu
doit trouver la même chose que les tronqués \inserted{au cran 2}, dans la catégorie \uline{homotopic}
des H-espaces  -  i.e. ce qu'on déduit des H-espaces par tuage des group
d'homotopie en dimensions \hand{$>$}2 .

\margin{\hand{$\|$} (double vertical line in the left margin, from here to ``lacets'')}(Demander à Quillen relations exactes entre \xout{H}objets-groupes et objet
\inserted[typed]{\xout{\unsure{pointésx}}}connexes de la catégorie homotopique\inserted{ponctuée}: forment-ils des catégories équival
tes par le foncteur espace des lacets\hand{?}).
\margin{\unsure{non, on a} \unsure{un foncteur} \unsure{pl. fidèle} (?) \illegible{3} \unsure{négatifs} \illegible{1}}

Pour les \uline{vrais} groupes \inserted{abéliens}, de (2-\unsure{\over{C}{C}}at), l'approche de Quillen consiste à
les identifier \unsure{à'abord|d'abord} à des groupes abéliens de la catégorie (Bicat),
donc à des grpupes abéliens bi-semiesimpliciaux particuliers (conditions
sur chaque ligne et chaque colonne). Appliquant la version bi-semisimpli
ale du tapis de Dold, on trouve que \xout{\illegible{1}} la catégorie de ces objets
est\xout{h} équivalente à la catégorie des \xout{objets} bicomplexes \inserted{de chaines} de groupes  abéli
qui sont nuls en dehors du carré  \xout{D}$0 \leq i,j \leq 1$.  \xout{Quillen dit que} La \deleted{colonn}\inserted{cond.}
que la première ligne corresponde à un \xout{ensemble} groupe semisimplicial
trivial signifie alors qu'elle n'a que le terme de degré zéro, i.e. que
le bicomplexe est réduit à

\diagram{
\begin{tikzcd}
C_{oo} & 0 \arrow[l] & \hand{$0$} \arrow[l, "\hand{}"'] \\
\hand{C}C_{1o} \arrow[u] & C_{11} \arrow[l] \arrow[u, "\hand{}"] & \hand{$0$} \arrow[l, hook', "\hand{}"']
\end{tikzcd}
}

\margin{\hand{9}}
